\subsection{覆叠的次数}

设 B 为拓扑空间，E 为 B 的覆叠，并用 p 表示其投影。设 C 为基数 \(\operatorname{Card}(\overline{p}^{1}(b))\) （当 b 取遍 B 时）所成的集。映射 \(b \mapsto \operatorname{Card}(\overline{p}^{1}(b))\) 是 B 到 C 中的局部常值映射。若 B 非空且映射 \(b \mapsto \operatorname{Card}(\overline{p}^{1}(b))\) 是常值的，则称覆叠 E 有次数。诸基数 \(\operatorname{Card}(\overline{p}^{1}(b))\) （对 \(b \in B\)）的公共值这时称为覆叠 E 的次数，记作 \(\deg(E, p)\)，或在不致与映射 p 相混淆时，简记作 \([E:B]\) 。

若 B 非空，则以 B 为基、以 F 为纤维型的平凡覆叠有次数，等于 Card(F)。

若 B 是连通的，则函数 \(b \mapsto \text{Card}(\overline{p}^{1}(b))\) 是常值的。于是：

命题 4. --- 非空连通空间上的每个覆叠都有次数。

设 G 为拓扑空间，\((\mathrm{F}, g)\) 为 G 的覆叠，\((\mathrm{E}, f)\) 为 F 的覆叠；设这些覆叠都有次数。若 G-空间 \((\mathrm{E}, g \circ f)\) 是覆叠，则它有次数，且 \(\deg(\mathrm{E}, g \circ f) = \deg(\mathrm{F}, g) \deg(\mathrm{E}, f)\)，它也可写成

\[
[ \mathrm{E}: \mathrm{G} ] = [ \mathrm{E}: \mathrm{F} ] [ \mathrm{F}: \mathrm{G} ].
\]

事实上，若 z 是 G 的一点，则映射

\[
f _ {g (z)} ^ {- 1}: \stackrel {{- 1}} {{f}} (\stackrel {{- 1}} {{g}} (z)) \to \stackrel {{- 1}} {{g}} (z)
\]

的诸纤维都有基数 \(\deg(\mathrm{E}, f)\)，而 \(\overline{g}^1(z)\) 的基数为 \(\deg(\mathrm{F}, g)\)。于是断言由牧羊人原理 (E, III, p. 41, prop. 9) 得出。

设 B 与 B' 为拓扑空间，(E, p) 为 B 的覆叠，(E', p') 为 B' 的覆叠。设这些覆叠都有次数；则 B × B' 的覆叠 (E × E', p × p') (I, p. 71, prop. 2) 有次数，且：

\[
\deg (\mathrm{E} \times \mathrm{E} ^ {\prime}, p \times p ^ {\prime}) = \deg (\mathrm{E}, p) \deg (\mathrm{E} ^ {\prime}, p ^ {\prime}).
\]

事实上，对每个偶对 \((b,b^{\prime})\in\mathrm{B}\times\mathrm{B}^{\prime}\)，纤维 \((\mathrm{E}\times\mathrm{E}^{\prime})_{(b,b^{\prime})}\) 就是积 \(E_{b}\times E_{b^{\prime}}^{\prime}\) 。

设 B 与 B' 为拓扑空间，(E, p) 为 B 的覆叠，(E', p') 为 B' 的覆叠，并设

\[
\begin{array}{c} \mathrm{E} ^ {\prime} \xrightarrow {f ^ {\prime}} \mathrm{E} \\ \Biggl \downarrow p ^ {\prime} \\ \mathrm{B} ^ {\prime} \xrightarrow {f} \mathrm{B} \end{array}
\]

为一个笛卡尔方块。若 B′ 非空且覆叠 (E, p) 有次数，则覆叠 (E′, p′) 有次数，等于 E 的次数 (I, p. 10, cor. of prop. 4)。

